\lyseis{3.4}{Τυποποιημένη μορφή μεγάλων αριθμών}
\sol{1} α) $7{,}2\cdot10^4$ \quad β) $1{,}25\cdot10^6$ \quad γ) $9{,}03\cdot10^8$ \quad δ) $6\cdot10^9$ \quad ε) $4{,}87\cdot10^4$ \quad στ) $1{,}5\cdot10^{13}$.
\sol{2} α) $270\,000$ \quad β) $5\,060$ \quad γ) $19\,000\,000$ \quad δ) $80\,000$ \quad ε) $3\,125\,000$ \quad στ) $4\,008\,000\,000$.
\sol{3} α) Ναι. β) Όχι: ο $52$ δεν είναι μικρότερος από $10$. γ) Όχι: ο $0{,}52$ είναι μικρότερος από $1$. δ) Όχι: ο $10$ δεν είναι μικρότερος από $10$. ε) Ναι.
\sol{4} α) $24\cdot10^5=2{,}4\cdot10^6$ \quad β) $0{,}35\cdot10^7=3{,}5\cdot10^6$ \quad γ) $150\cdot10^3=1{,}5\cdot10^5$ \quad δ) $0{,}008\cdot10^9=8\cdot10^6$.
\sol{5} α) $>$ (εκθέτες $5>4$) \quad β) $>$ (ίδιος εκθέτης, $3{,}5>2{,}9$) \quad γ) $>$ ($6>5{,}99$) \quad δ) $<$ (εκθέτες $3<5$).
\sol{6} $7{,}9\cdot10^5<9\cdot10^5<3{,}2\cdot10^6<4{,}5\cdot10^6<1{,}1\cdot10^7$.
\sol{7} α) $45\,000\,000$ \quad β) $25\,000\,000\,000$. Το γινόμενο είναι $36\,000\,000\,000\,000$ και η οθόνη δείχνει $3.6\mathrm{E}13$.
\sol{8} $3{,}844\cdot10^5$ χλμ. και $1{,}496\cdot10^8$ χλμ. Ο λόγος είναι περίπου $\frac{1{,}496\cdot10^8}{3{,}844\cdot10^5}\approx389$, δηλαδή περίπου $390$ φορές μακρύτερα.
\sol{9} $9{,}461\cdot10^{12}$ χλμ. Τα $4{,}3$ έτη φωτός είναι $4{,}3\cdot9{,}461\cdot10^{12}=40{,}68\cdot10^{12}=4{,}07\cdot10^{13}$ χλμ.
\sol{10} $\frac{1{,}496\cdot10^8}{3\cdot10^5}\approx499$ δευτερόλεπτα, δηλαδή $499:60\approx8{,}3$ λεπτά.
\sol{11} $8\cdot10^9$ και $1{,}04\cdot10^7$. Ο λόγος είναι $\frac{8\cdot10^9}{1{,}04\cdot10^7}\approx769$, περίπου $770$ φορές.
\sol{12} $5\,976\,000\,000\,000\,000\,000\,000\,000$ κιλά. Ο λόγος $\frac{1{,}983\cdot10^{30}}{5{,}976\cdot10^{24}}\approx3{,}3\cdot10^5$, δηλαδή περίπου $330\,000$ φορές.
\sol{13} $\frac{1{,}28\cdot10^{11}}{4\cdot10^6}=0{,}32\cdot10^5=3{,}2\cdot10^4=32\,000$ φωτογραφίες.
\sol{14} Επειδή ο πρώτος παράγοντας $10$ δεν είναι μικρότερος από $10$. Γράφεται $10\cdot10^5=10^6=1\cdot10^6$.
\sol{15} Ο $4{,}56\cdot10^{10}$ έχει $11$ ψηφία ($\nu+1$, όπου $\nu=10$ ο εκθέτης). Ο $1\cdot10^{\nu}$ έχει $\nu+1$ ψηφία: το $1$ και $\nu$ μηδενικά.
\sol{16} $(3\cdot10^4)\cdot(2\cdot10^5)=6\cdot10^9$. $(4\cdot10^3)\cdot(5\cdot10^6)=20\cdot10^9=2\cdot10^{10}$. Κανόνας: πολλαπλασιάζω τους $\alpha$, προσθέτω τους εκθέτες και, αν ο νέος $\alpha$ βγει ίσος ή μεγαλύτερος από $10$, ξαναγράφω σε τυποποιημένη μορφή.
\sol{17} Όχι. Ο $9{,}9\cdot10^3=9\,900$ και ο $1\cdot10^4=10\,000$, άρα $9{,}9\cdot10^3<1\cdot10^4$. Συγκρίνουμε πρώτα τους εκθέτες.
\sol{18} $\frac{10^{15}}{3{,}15\cdot10^7}\approx3{,}17\cdot10^7$ έτη, δηλαδή περίπου $32$ εκατομμύρια χρόνια.
\sol{19} Ελεύθερη απάντηση. Π.χ. μάζα του Ήλιου $1{,}989\cdot10^{30}$ κιλά, ακτίνα της Γης $6{,}371\cdot10^3$ χλμ. (οι τιμές διαφέρουν λίγο ανά πηγή).
\sol{20} Ελεύθερη απάντηση. Έλεγξε ότι ο $\alpha$ είναι τουλάχιστον $1$ και λιγότερο από $10$ και ότι ο εκθέτης είναι ίσος με το πλήθος των θέσεων της υποδιαστολής.
